% ----------------------------------------------------------------------
%  2.2.2. Движение дислокаций
% ----------------------------------------------------------------------
\subsection{Движение дислокаций}\label{ssec:disl-motion}

\begin{orientir}
Дислокация движется двумя принципиально разными способами.
\textbf{Скольжение}~— перемещение в плоскости скольжения; атомы лишь
слегка сдвигаются в ядре, переноса вещества нет, поэтому оно возможно
при любой температуре. \textbf{Переползание}~— перемещение краевой
дислокации поперёк плоскости скольжения; для него нужно подвести или
отвести атомы, т.\,е. нужна диффузия. В конце раздела мы свяжем движение
дислокаций с макроскопической пластической деформацией (формула
Орована).
\end{orientir}

\begin{opredelenie}
\term{Консервативное движение} (\term{скольжение})~— движение дислокации,
не сопровождающееся переносом массы. \term{Неконсервативное движение}
(\term{переползание})~— движение, требующее переноса массы (притока или
оттока атомов).
\end{opredelenie}

\subsubsection{Скольжение краевой дислокации}

\paragraph{Плоскость скольжения}
Краевая дислокация скользит в направлении своего вектора Бюргерса.
Плоскость, содержащая взаимно перпендикулярные векторы $\vb b$ и $\vb t$,
единственна; это \term{плоскость скольжения}. Сойти с неё скольжением
краевая дислокация не может.

\paragraph{Атомный механизм}
Чтобы дислокация сместилась на одно межатомное расстояние, вовсе не нужно
сдвигать всю верхнюю половину кристалла. Достаточно, чтобы немного
сместились атомы вблизи ядра (рис.~\ref{fig:edge-glide}). Одна
вертикальная атомная плоскость рядом с экстраплоскостью «разрывается»:
её нижняя половина объединяется с экстраплоскостью в целую плоскость, а
верхняя становится новой экстраплоскостью. Каждый атом при этом
смещается лишь на долю межатомного расстояния.

\begin{figure}[ht]
\centering
\begin{subfigure}[b]{0.48\textwidth}\centering
\begin{tikzpicture}[scale=0.45]
  \begin{scope}\clip (-4.6,-3.6) rectangle (4.6,3.6);
  \EdgeDislLattice{6}{4}{4}\end{scope}
  \dislT{(0,-0.05)}{0}
  \draw[arr, fmRed, line width=1pt] (-0.4,4.0) -- (1.4,4.0) node[right, slbl]{$\tau$};
  \draw[arr, fmRed, line width=1pt] (1.4,-4.0) -- (-0.4,-4.0) node[left, slbl]{$\tau$};
\end{tikzpicture}
\caption{}
\end{subfigure}\hfill
\begin{subfigure}[b]{0.48\textwidth}\centering
\begin{tikzpicture}[scale=0.45]
  \begin{scope}\clip (-4.6,-3.6) rectangle (4.6,3.6);
  \EdgeDislLattice[xshift=1cm]{6}{4}{4}\end{scope}
  \dislT{(1,-0.05)}{0}
  \draw[fmGray, dashed] (0,-0.2) -- (0,3.6);
\end{tikzpicture}
\caption{}
\end{subfigure}
\caption{Скольжение краевой дислокации на одно межатомное расстояние
вправо: \textit{а}~— исходное положение; \textit{б}~— после смещения.
Новая экстраплоскость (оранжевая)~— это верхняя половина бывшей соседней
плоскости, а бывшая экстраплоскость (пунктир) соединилась с нижней
половиной той плоскости и стала целой. Атомы сместились лишь на доли
межатомного расстояния}
\label{fig:edge-glide}
\end{figure}

\begin{glavnoe}
Скольжение не требует диффузии: атомы в ядре смещаются на малые
расстояния согласованно, под действием напряжения. Поэтому дислокации
скользят \emph{при сколь угодно низких температурах}.
\end{glavnoe}

Приложенное к кристаллу напряжение создаёт силу, действующую на линию
дислокации. Её составляющая в плоскости скольжения, приходящаяся на
единицу длины линии, равна
\begin{equation}\label{eq:F-glide}
F=\tau b ,
\end{equation}
где $\tau$~— касательное напряжение в плоскости скольжения в направлении
$\vb b$ (общий вывод~— формула Пича--Кёлера, п.~\ref{ssec:disl-elastic}).
Это согласуется с опытным \emph{законом Шмида}: скольжение начинается,
когда \emph{приведённое} касательное напряжение в системе скольжения
достигает критического значения (п.~1.6.1).

\subsubsection{Напряжение Пайерлса}

Даже в совершенном кристалле дислокации приходится преодолевать
сопротивление самой решётки. Пусть дислокация параллельна
плотноупакованному ряду (рис.~\ref{fig:peierls},\,а). Переходя из одного
положения равновесия в соседнее, её ядро проходит через промежуточные
конфигурации с несколько большей энергией. Энергия дислокации~—
периодическая функция её положения (рис.~\ref{fig:peierls},\,б): это
\term{рельеф Пайерлса}. Чтобы перевалить через барьер при $T=0$, нужно
приложить напряжение $\tau_\text{П}$~— \term{напряжение Пайерлса}
(Пайерлса--Набарро).

\begin{figure}[ht]
\centering
\begin{subfigure}[b]{0.4\textwidth}\centering
\begin{tikzpicture}[scale=0.62]
  \foreach \y in {0,1,2}{\fill[fmBlue!8] (-0.3,{\y-0.25}) rectangle (5.3,{\y+0.25});}
  \foreach \y in {0,1,2}{\draw[fmBlue!40, dashed] (-0.3,\y) -- (5.3,\y);}
  \draw[disl] (-0.3,0) -- (5.3,0);
  \draw[arr, fmGray] (2.5,0.15) -- (2.5,0.85);
  \node[slbl, above] at (2.5,2.3) {«долины» рельефа};
\end{tikzpicture}
\caption{}
\end{subfigure}\hfill
\begin{subfigure}[b]{0.56\textwidth}\centering
\begin{tikzpicture}
\begin{axis}[width=\linewidth, height=4cm, xmin=0, xmax=3, ymin=0, ymax=1.3,
  xtick={0,1,2,3}, xticklabels={$0$,$b$,$2b$,$3b$}, ytick=\empty,
  xlabel={положение дислокации}, ylabel={$E$}]
  \addplot[fmBlue, very thick, domain=0:3, samples=150] {0.5*(1-cos(deg(2*pi*x)))+0.1};
  \draw[dimline] (axis cs:1.5,0.1) -- node[right, font=\footnotesize]{$\Delta E_{\text{П}}$} (axis cs:1.5,1.1);
\end{axis}
\end{tikzpicture}
\caption{}
\end{subfigure}
\caption{Рельеф Пайерлса: \textit{а}~— дислокация лежит в одной из
«долин» потенциальной энергии, параллельных плотноупакованным рядам;
\textit{б}~— энергия дислокации (на единицу длины) периодически
зависит от её положения}
\label{fig:peierls}
\end{figure}

Оценка в модели Пайерлса--Набарро (области над и под плоскостью
скольжения~— упругие сплошные среды, связанные синусоидальным
взаимодействием) даёт
\begin{keyeq}\label{eq:peierls}
\tau_\text{П}\approx\frac{2G}{K}\Ex{-\frac{2\pi a}{Kb}}=\frac{2G}{K}\Ex{-\frac{2\pi w}{b}},
\qquad w=\frac aK ,
\end{keyeq}
где $a$~— расстояние между плоскостями скольжения, $b$~— межатомное
расстояние в направлении скольжения, $K=1$ для винтовой и $K=1-\nu$ для
краевой дислокации ($\nu$~— коэффициент Пуассона). Энергетический барьер
на единицу длины $\Delta E_\text{П}\approx\frac{Gb^2}{2\pi K}\exp(-2\pi
a/Kb)$.

\paragraph{Ширина дислокации}
Величина $w$~— \term{ширина дислокации}: ширина области в плоскости
скольжения, где смещения атомов больше половины максимального. Резкой
границы ядра нет, поэтому понятие условное; в металлах ширина составляет
от двух до десяти межатомных расстояний. Почему $\tau_\text{П}$ так
сильно (экспоненциально) зависит от ширины?

\begin{smysl}[Дислокация как нониус]
Краевая дислокация похожа на нониус штангенциркуля: $n$ атомных
плоскостей с одной стороны плоскости скольжения противостоят $n+1$
плоскостям с другой. При перемещении дислокации на $b$ каждая из
плоскостей сдвигается лишь на $b/n$. Чем шире дислокация (больше $n$),
тем меньше смещение каждого атома и тем меньшая сила нужна. Широкая
дислокация почти «не чувствует» дискретности решётки.
\end{smysl}

\paragraph{Какие плоскости лучше}
Из~\eqref{eq:peierls}: $\tau_\text{П}$ мало, когда $a/b$ велико~— т.\,е.
для \emph{плотноупакованных} плоскостей (они дальше всего отстоят друг от
друга) и \emph{коротких} векторов Бюргерса (направления плотной упаковки).
Именно поэтому скольжение идёт по плотноупакованным плоскостям в
плотноупакованных направлениях~— по \term{системам скольжения} (п.~1.6.1).
Например, для плоскости $\{111\}$ ГЦК-решётки $a=b\sqrt{2/3}$ и для
краевой дислокации $\tau_\text{П}\sim10^{-3}G$.

\begin{vnimanie}
Модель Пайерлса--Набарро очень приближённая и даёт лишь порядок величины
(и то для простых решёток): из-за экспоненты небольшая ошибка в $a/b$ или
$w$ меняет результат на порядки. Для металлов реальные значения обычно
ещё ниже. Надёжные оценки требуют расчёта атомной структуры ядра.
\end{vnimanie}

\paragraph{Когда напряжение Пайерлса велико}
Оно мало, если связи в ядре металлические: такие связи ненаправленные и
мало меняются при перестройке ядра. Оно велико, если связи ионные
(одноимённые ионы при скольжении проходят близко друг от друга) или
ковалентные (направленные связи приходится рвать), а также в сложных
структурах с несколькими атомами в ячейке. Отсюда грубая классификация
кристаллов (табл.~\ref{tab:peierls-classes}).

\begin{table}[ht]
\centering
\caption{Классификация кристаллов по лёгкости скольжения}
\label{tab:peierls-classes}
\small
\begin{tabularx}{\textwidth}{@{}>{\raggedright\arraybackslash}p{38mm}>{\raggedright\arraybackslash}X@{}}
\toprule
Скольжение & Примеры \\
\midrule
лёгкое при всех температурах & «обычные» металлы и сплавы (Na, Cu, Al и т.\,п.);
  слоистые структуры вдоль слоёв (графит) \\
лёгкое при комнатной и высокой температуре, трудное при очень низкой &
  ОЦК переходные металлы с долей ковалентной связи (Cr, Fe, W); ионные кристаллы \\
только при высоких температурах & ковалентные кристаллы (алмаз, Si); слоистые
  структуры поперёк слоёв; сложные структуры ($\alpha$-Mn, интерметаллиды) \\
\bottomrule
\end{tabularx}
\end{table}

\subsubsection{Перегибы: как дислокация обходит барьер Пайерлса}

Дислокации не обязательно преодолевать барьер Пайерлса сразу по всей
длине. Её участок может «перескочить» в соседнюю долину рельефа, так что
на линии образуется \term{двойной перегиб} (рис.~\ref{fig:kinks}). Дальше
перегибы легко расходятся в стороны вдоль долины~— для этого нужно
напряжение много меньше $\tau_\text{П}$,~— и вся дислокация участок за
участком переходит в соседнюю долину.

\begin{figure}[ht]
\centering
\begin{tikzpicture}[scale=0.8]
  \foreach \y in {0,1,2}{\fill[fmBlue!8] (-0.3,{\y-0.25}) rectangle (10.3,{\y+0.25});
                         \draw[fmBlue!40, dashed] (-0.3,\y) -- (10.3,\y);}
  \draw[disl, rounded corners=6pt] (-0.3,0) -- (3.4,0) -- (4.0,1) -- (6.0,1) -- (6.6,0) -- (10.3,0);
  \draw[arr, fmRed] (3.6,0.5) -- (2.6,0.5);
  \draw[arr, fmRed] (6.4,0.5) -- (7.4,0.5);
  \node[slbl, below] at (3.7,-0.25) {перегиб};
  \node[slbl, below] at (6.3,-0.25) {перегиб};
  \draw[arr, fmGray] (5,2.1) node[above, slbl]{$\tau$} -- (5,1.4);
\end{tikzpicture}
\caption{Двойной перегиб: участок дислокации перешёл в соседнюю долину
рельефа Пайерлса. Перегибы расходятся вдоль линии (красные стрелки), и
дислокация постепенно целиком смещается на один период}
\label{fig:kinks}
\end{figure}

Зарождение двойного перегиба требует тепловой флуктуации, поэтому скорость
дислокации сильно зависит от температуры:
\begin{equation}\label{eq:kink-vel}
V=A\Ex{-\frac{E_{\text{пер}}}{kT}},
\end{equation}
где $E_{\text{пер}}$~— энергия образования двойного перегиба.

\begin{glavnoe}
Основная идея теории дислокаций~— \emph{неодновременность} сдвига:
кристалл сдвигается не целиком, а по мере пробега дислокации. Та же идея
повторяется на следующем уровне: и сама дислокация перемещается не
целиком, а по частям~— за счёт движения перегибов.
\end{glavnoe}

\subsubsection{Переползание краевой дислокации}

Пусть теперь краевая дислокация смещается \emph{перпендикулярно}
плоскости скольжения. Это значит, что экстраплоскость должна стать
короче или длиннее (рис.~\ref{fig:climb}).

\begin{figure}[ht]
\centering
\begin{subfigure}[b]{0.48\textwidth}\centering
\begin{tikzpicture}[scale=0.45]
  \begin{scope}\clip (-4.6,-3.6) rectangle (4.6,3.6);
  \EdgeDislLattice{6}{4}{4}\end{scope}
  \dislT{(0,-0.05)}{0}
  \node[vac=3.6mm] at (U1_0) {};
  \draw[arr, fmRed, line width=1pt] (U0_0) ++(0.2,-0.25) .. controls +(0.4,-0.35) .. ($(U1_0)+(-0.15,-0.2)$);
\end{tikzpicture}
\caption{}
\end{subfigure}\hfill
\begin{subfigure}[b]{0.48\textwidth}\centering
\begin{tikzpicture}[scale=0.45]
  \begin{scope}\clip (-4.6,-3.6) rectangle (4.6,3.6);
  \EdgeDislLattice[yshift=1cm]{6}{4}{5}\end{scope}
  \dislT{(0,0.95)}{0}
  \draw[fmGray, dashed] (-4.6,0) -- (4.6,0);
  \draw[fmGray, dashed] (-4.6,1) -- (4.6,1);
\end{tikzpicture}
\caption{}
\end{subfigure}
\caption{Переползание краевой дислокации: \textit{а}~— к кромке
экстраплоскости подошла вакансия, и крайний атом перескакивает в неё;
\textit{б}~— после ухода всего ряда атомов (вдоль линии дислокации)
экстраплоскость стала на один атом короче, и дислокация поднялась в
соседнюю плоскость скольжения}
\label{fig:climb}
\end{figure}

\begin{itemize}
\item \term{Положительное переползание} (экстраплоскость укорачивается,
дислокация $\perp$ поднимается): атомы с кромки экстраплоскости уходят
в кристалл. Либо к кромке подходят вакансии и атомы перескакивают в
них, либо атомы с кромки переходят в междоузлия и диффундируют прочь.
Первый путь гораздо вероятнее: избыточные вакансии в металле возникают
легко, а энергия образования межузельного атома велика.
\item \term{Отрицательное переползание} (экстраплоскость
достраивается): к кромке присоединяются межузельные атомы, либо к ней
присоединяются соседние атомы из узлов, порождая вакансии, которые
уходят вглубь (этот вариант маловероятен).
\end{itemize}

\begin{glavnoe}
Переползание~— диффузионный, термически активируемый процесс. Его
скорость сильно зависит от температуры, и при низких температурах
переползания практически нет. Именно этим переползание отличается от
скольжения.
\end{glavnoe}

\paragraph{Пороги}
Вакансии подходят к дислокации поодиночке, поэтому дислокация переползает
не сразу по всей длине, а по частям. На линии образуются ступеньки
высотой в межатомное расстояние~— \term{пороги}
(рис.~\ref{fig:jogs}). Порог~— место, где удобнее всего присоединяются и
отрываются атомы: присоединение вакансии к порогу просто сдвигает его
вдоль линии. Переползание~— это зарождение порогов и их движение вдоль
дислокации.

\begin{figure}[ht]
\centering
\begin{tikzpicture}[cubeproj, scale=1.1]
  \foreach \z/\c in {0/fmBlue!7, 0.6/fmBlue!12}{
    \filldraw[fill=\c, draw=fmBlue!40] (0,0,\z) -- (6,0,\z) -- (6,2.2,\z) -- (0,2.2,\z) -- cycle;}
  \draw[disl, line width=1.6pt, line join=round] (0,1,0) -- (1.8,1,0) -- (1.8,1,0.6) -- (3.6,1,0.6) -- (3.6,1,0) -- (4.5,1,0) -- (4.5,1,0.6) -- (6,1,0.6);
  \node[slbl, above] at (1.8,1,0.6) {порог};
  \node[slbl, above] at (4.5,1,0.6) {порог};
  \node[slbl, right] at (6,2.2,0.6) {плоскость $P'$};
  \node[slbl, right] at (6,2.2,0) {плоскость $P$};
\end{tikzpicture}
\caption{Пороги на краевой дислокации: участки линии лежат в соседних
параллельных плоскостях скольжения $P$ и $P'$, а пороги соединяют их}
\label{fig:jogs}
\end{figure}

В равновесии число порогов на единице длины дислокации
\[
N=n_0\Ex{-\frac{E_{\text{пор}}}{kT}},
\]
где $n_0$~— число атомов на единице длины линии, $E_{\text{пор}}\sim1\eV$~—
энергия образования порога. Энергия активации переползания складывается
из энергии образования порога и энергии активации самодиффузии:
$E=E_{\text{пор}}+E_D$. Если на дислокации уже много порогов (например,
после пересечения с другими дислокациями), то $E=E_D$.

\paragraph{Сколько вещества нужно перенести}
Общий подход годится для дислокации любого типа. Пусть участок
дислокации, двигаясь, заметает площадку $\d S$ с единичной нормалью
$\vb n$ (знак $\vb n$ задаётся правилом $\vb n=\vb t\times\vb d$, где
$\vb d$~— направление движения). Материал над площадкой смещается на
$\vb b$ относительно материала под ней, и объём, который надо
добавить или убрать, равен
\begin{equation}\label{eq:climb-volume}
\d V=(\vb b\cdot\vb n)\,\d S .
\end{equation}
\begin{itemize}
\item Если $\vb b\cdot\vb n=0$ (плоскость движения содержит $\vb b$), то
переноса массы нет~— это скольжение.
\item Если $\vb b\cdot\vb n\neq0$~— переползание: при $\vb b\cdot\vb n>0$
образуется пустота (атомы должны подойти), при $\vb b\cdot\vb n<0$~—
избыток материала (атомы должны уйти).
\end{itemize}
Для винтовой дислокации $\vb b\parallel\vb t$, и любая плоскость,
содержащая линию, содержит и $\vb b$. Поэтому \vazhno{винтовая
дислокация не переползает}~— она может только скользить.

\begin{podrobnee}[Сила, необходимая для переползания, и «осмотическая» сила]
Пусть линия дислокации составляет угол $\varphi$ с вектором $\vb b$ и
переползает, заметая площадь $\d S$. Объём оставленной пустоты (или
лишнего материала) $\d V=b\sin\varphi\,\d S$, что эквивалентно
$\d N=\d V/b^3=\d S\sin\varphi/b^2$ точечным дефектам. Если на их создание
тратится энергия $E^{\mathrm F}$, то на единицу длины дислокации должна
действовать сила
\[
F_m=E^{\mathrm F}\frac{\d N}{\d S}=\frac{E^{\mathrm F}}{b^2}\sin\varphi .
\]
Для дислокаций, не близких к винтовым, эта сила огромна~— порядка
теоретического предела упругости. Значит, \emph{без пересыщения
точечными дефектами} переползание идёт лишь в особых условиях.

Пусть теперь кристалл пересыщен вакансиями: их концентрация $C$ больше
равновесной $C_0$. Тогда поглощение вакансий дислокацией выгодно, и
пересыщение действует на неё как сила, перпендикулярная плоскости
скольжения (\emph{осмотическая}, или \emph{химическая} сила):
\[
F_S=\frac{kT\sin\varphi}{b^2}\ln\frac{C}{C_0} .
\]
Например, после закалки от $T_1$ до $T_2$ пересыщение
$C/C_0=\exp\bigl[(E^{\mathrm F}/k)(1/T_2-1/T_1)\bigr]$, и
\[
F_S=F_m\,\frac{T_1-T_2}{T_1}
\]
~— это значительная доля $F_m$. Холодная деформация создаёт концентрации
вакансий и межузельных атомов того же порядка, что и нагрев почти до
плавления. Поэтому после закалки или деформации для переползания
достаточно напряжений, много меньших, чем в кристалле без пересыщения.
\end{podrobnee}

\subsubsection{Скольжение винтовой дислокации}

Винтовая дислокация тоже скользит, но иначе, чем краевая
(рис.~\ref{fig:screw-glide}).

\begin{figure}[ht]
\centering
\begin{subfigure}[b]{0.48\textwidth}\centering
\begin{tikzpicture}[scale=0.42]
  \begin{scope}\clip (-5.5,-0.6) rectangle (5.5,6.6);\ScrewDislPlan{6}{6}\end{scope}
  \draw[disl, line width=1.2pt] (0,-0.6) -- (0,6.6);
\end{tikzpicture}
\caption{}
\end{subfigure}\hfill
\begin{subfigure}[b]{0.48\textwidth}\centering
\begin{tikzpicture}[scale=0.42]
  \begin{scope}\clip (-5.5,-0.6) rectangle (5.5,6.6);\ScrewDislPlan[xshift=1cm]{7}{7}\end{scope}
  \draw[disl, line width=1.2pt] (1,-0.6) -- (1,6.6);
  \draw[fmGray, dashed] (0,-0.6) -- (0,6.6);
\end{tikzpicture}
\caption{}
\end{subfigure}
\caption{Скольжение винтовой дислокации на одно межатомное расстояние
вправо (вид на плоскость скольжения, обозначения как на
рис.~\ref{fig:screw-block},\,б). Атомы ядра смещаются \emph{вдоль} линии
дислокации (вверх и вниз на рисунке), а сама линия перемещается
\emph{поперёк}~— вправо}
\label{fig:screw-glide}
\end{figure}

Отметим три важных обстоятельства:
\begin{enumerate}[label=\arabic*)]
\item дислокация смещается на период, хотя атомы двигаются только внутри
ядра;
\item каждый атом ядра смещается лишь на долю периода (как и у краевой
дислокации);
\item атомы смещаются \emph{параллельно} действующим на них силам и
вектору $\vb b$ (вдоль линии), а сама дислокация движется
\emph{перпендикулярно} им. Этим винтовая дислокация отличается от
краевой, которая движется вдоль $\vb b$.
\end{enumerate}

Когда винтовая дислокация пройдёт через весь кристалл, одна его часть
окажется сдвинутой относительно другой на $\vb b$~— так же, как после
пробега краевой. Под действием одного и того же напряжения винтовые
дислокации разного знака (правая и левая) движутся в противоположные
стороны.

\paragraph{Поперечное скольжение}
У винтовой дислокации нет единственной плоскости скольжения: любая
плоскость, проходящая через линию, содержит и $\vb b$. Поэтому, встретив
препятствие в плоскости $P$, винтовая дислокация может продолжить
скольжение в другой плоскости $R$, пересекающейся с $P$ по направлению
$\vb b$ (рис.~\ref{fig:cross-slip}). Это \term{поперечное скольжение}.
Обойдя препятствие, дислокация может перейти в плоскость $S$,
параллельную исходной~— это \term{двойное поперечное скольжение}.
Многократное повторение называют \term{множественным поперечным
скольжением}.

\begin{figure}[ht]
\centering
\begin{tikzpicture}[cubeproj, scale=1.05, line join=round]
  \def\Lx{6}
  \filldraw[fill=fmBlue!7, draw=fmBlue!45] (0,0,0) -- (\Lx,0,0) -- (\Lx,2,0) -- (0,2,0) -- cycle;
  \filldraw[fill=fmGreen!10, draw=fmGreen!55] (0,2,0) -- (\Lx,2,0) -- (\Lx,2.9,1.4) -- (0,2.9,1.4) -- cycle;
  \filldraw[fill=fmBlue!7, draw=fmBlue!45] (0,2.9,1.4) -- (\Lx,2.9,1.4) -- (\Lx,4.6,1.4) -- (0,4.6,1.4) -- cycle;
  \fill[black] (3,2.3,0) circle (2.5pt) node[below, slbl]{препятствие};
  \draw[disl, line width=1.7pt] (0,1.5,0) -- (1.6,1.5,0) -- (1.9,2,0) -- (2.15,2.45,0.7) -- (2.4,2.9,1.4)
       -- (2.5,3.6,1.4) -- (3.6,3.6,1.4) -- (3.75,2.9,1.4) -- (4.0,2.45,0.7) -- (4.25,2,0) -- (4.5,1.5,0) -- (\Lx,1.5,0);
  \draw[bvec] (0.4,0.4,0) -- (1.6,0.4,0) node[below right, slbl]{$\vb b$};
  \node[slbl, fmBlue] at (5.4,0.3,0) {$P$};
  \node[slbl, fmGreen] at (5.6,2.5,0.75) {$R$};
  \node[slbl, fmBlue] at (5.4,4.3,1.4) {$S$};
\end{tikzpicture}
\caption{Двойное поперечное скольжение винтовой дислокации: встретив
препятствие в плоскости $P$, участок дислокации переходит в
пересекающую плоскость $R$, а затем в плоскость $S$, параллельную $P$.
Все три плоскости содержат вектор $\vb b$}
\label{fig:cross-slip}
\end{figure}

\begin{glavnoe}
Краевая дислокация может сменить плоскость скольжения только
переползанием (с переносом массы, диффузионно). Винтовая~— поперечным
скольжением, без переноса массы.
\end{glavnoe}

\subsubsection{Смешанные дислокации}

Дислокация, вектор Бюргерса которой не параллелен и не перпендикулярен
линии, называется \term{смешанной}. В общем случае дислокация~— это
граница между частью плоскости, где сдвиг уже прошёл, и частью, где его
ещё нет; такая граница может быть искривлённой. Под действием
касательного напряжения каждый участок движется перпендикулярно своей
линии: краевые участки~— вдоль $\vb b$, винтовые~— поперёк $\vb b$, и зона
сдвига расширяется (вспомните петлю на рис.~\ref{fig:loop-types}). Когда
вся линия выйдет на поверхность (разные участки выходят в разное время),
верхняя часть кристалла окажется сдвинутой относительно нижней на $\vb b$.

\begin{glavnoe}
Выход любой дислокации на поверхность сдвигает одну часть кристалла
относительно другой на вектор Бюргерса этой дислокации.
\end{glavnoe}

\subsubsection{Пластическая деформация как движение дислокаций}

\paragraph{Один пробег~— одна ступенька}
Сдвинем верхнюю половину кубического кристалла размером $L$ по плоскости
$AA$ на вектор $\vb b$ (рис.~\ref{fig:disl-strain}). Это можно сделать
двумя способами:
\begin{enumerate}[label=\arabic*)]
\item сдвинуть всю верхнюю часть сразу~— нужно напряжение порядка
теоретической прочности;
\item создать на плоскости $AA$ краевую дислокацию с вектором $\vb b$ и
прогнать её через весь кристалл~— нужно напряжение порядка
$\tau_\text{П}$.
\end{enumerate}
Результат одинаков: сдвиг на $b$ и деформация $\varepsilon\approx b/L$.
Если прошли $n$ дислокаций, $\varepsilon\approx nb/L$.

\begin{figure}[ht]
\centering
\begin{tikzpicture}[scale=0.62, line join=round]
  \def\bb{0.45}
  \begin{scope}
    \filldraw[fill=fmBlue!12, draw=atomLine] (0,0) rectangle (3,3);
    \draw[dashed, fmGray] (-0.3,1.5) -- (3.3,1.5) node[right, slbl]{$A$};
    \node[lbl] at (1.5,-0.6) {а};
  \end{scope}
  \begin{scope}[shift={(5,0)}]
    \filldraw[fill=fmBlue!12, draw=atomLine] (0,0) -- (3,0) -- (3,1.5) -- (3+\bb,1.5) -- (3+\bb,3) -- (\bb,3) -- (\bb,1.5) -- (0,1.5) -- cycle;
    \draw[dimline] (3,1.2) -- node[below, slbl]{$b$} (3+\bb,1.2);
    \node[lbl] at (1.5,-0.6) {б};
  \end{scope}
  \begin{scope}[shift={(10,0)}]
    \filldraw[fill=fmBlue!12, draw=atomLine] (0,0) -- (3,0) -- (3,3) -- (\bb,3) -- (\bb,1.5) -- (0,1.5) -- cycle;
    \dislT[fmRed]{(0.9,1.5)}{0}
    \draw[arr, fmRed] (1.3,1.75) -- (2.3,1.75);
    \node[lbl] at (1.5,-0.6) {в};
  \end{scope}
  \begin{scope}[shift={(15,0)}]
    \filldraw[fill=fmBlue!12, draw=atomLine] (0,0) -- (3,0) -- (3,3) -- (\bb,3) -- (\bb,1.5) -- (0,1.5) -- cycle;
    \dislT[fmRed]{(2.1,1.5)}{0}
    \draw[dimline] (\bb,1.15) -- node[below, slbl]{$l$} (2.1,1.15);
    \node[lbl] at (1.5,-0.6) {г};
  \end{scope}
\end{tikzpicture}
\caption{Пластическая деформация движением дислокации: \textit{а}~—
исходный кристалл; \textit{б}~— верхняя половина сдвинута по плоскости
$A$ на $b$; \textit{в}~— на левой грани зародилась краевая дислокация;
\textit{г}~— дислокация прошла путь $l$ (незавершённый сдвиг). Когда она
выйдет на правую грань, получится состояние~\textit{б}}
\label{fig:disl-strain}
\end{figure}

Дислокация может пройти лишь часть пути $l<L$ и остановиться у
препятствия; $l$~— \term{длина свободного пробега}, а сдвиг называют
\emph{незавершённым} (само понятие дислокации возникло при наблюдении
незавершённых сдвигов; термин заимствован из геологии, где так называли
незавершённые сдвиги горных пород). Тогда удлинение пропорционально
пройденной доле: $\Delta L=bl/L$. Для $n$ дислокаций в кристалле сечением
$L^2$:
\[
\varepsilon=\frac{n}{L^2}\,b\,l .
\]
Число дислокаций, пересекающих единицу площади, $n/L^2=\rho_\text{п}$~—
\term{плотность подвижных дислокаций} (размерность см$^{-2}$; то же, что
суммарная длина дислокаций в единице объёма).

\begin{keyeq}\label{eq:orowan}
\varepsilon=\rho_\text{п}\,b\,l,\qquad \dot\varepsilon=\rho_\text{п}\,b\,v
\quad\text{(формула Орована)} .
\end{keyeq}
Вторая формула получается дифференцированием первой по времени: полностью
$\dot\varepsilon=\rho_\text{п}bv+\dot\rho_\text{п}bl$, где $v$~— скорость
дислокаций, но при небольших напряжениях изменение плотности подвижных
дислокаций $\dot\rho_\text{п}$ обычно мало, и вторым членом пренебрегают.

\begin{smysl}
Формула Орована связывает микроскопику (сколько дислокаций, с каким
вектором Бюргерса и как быстро движутся) с макроскопикой (скорость
деформации образца). Это основное уравнение дислокационной теории
пластичности.
\end{smysl}

\paragraph{Скорость дислокаций}
Прямые наблюдения показывают, что скорость дислокации зависит от
напряжения очень круто и не зависит от длительности импульса нагрузки
(рис.~\ref{fig:v-tau}). В интервале скоростей $10^{-7}$--$10^{-1}$~см/с
эмпирически
\begin{equation}\label{eq:v-tau}
v=v_0\left(\frac{\tau}{\tau_0}\right)^{m},
\end{equation}
где $\tau_0$~— напряжение, при котором $v_0=1$~см/с, $\tau$~—
касательное напряжение в плоскости скольжения, $m$~— константа
материала (часто 10--40). При больших напряжениях рост замедляется, и
скорость асимптотически приближается к скорости звука $c$.

\begin{figure}[ht]
\centering
\begin{tikzpicture}
\begin{axis}[width=0.6\textwidth, height=5.4cm, xmode=log, ymode=log,
  xmin=0.3, xmax=20, ymin=1e-8, ymax=3,
  xlabel={касательное напряжение $\tau$ (отн. ед.)}, ylabel={$v/c$},
  ytick={1e-8,1e-6,1e-4,1e-2,1}]
  \addplot[fmBlue, very thick, domain=0.3:20, samples=150] {(x/2)^12/(1+(x/2)^12)^(1)};
  \draw[dashed, fmGray] (axis cs:0.3,1) -- (axis cs:20,1) node[pos=0.2, above, font=\footnotesize]{скорость звука};
  \node[font=\footnotesize, rotate=58, fmBlue] at (axis cs:0.85,2e-4) {$v\propto\tau^{m}$};
\end{axis}
\end{tikzpicture}
\caption{Зависимость скорости дислокаций от напряжения (схема, логарифмические
оси). На крутом участке $v\propto\tau^m$; при больших напряжениях скорость
стремится к скорости звука}
\label{fig:v-tau}
\end{figure}

\begin{primer}[Хватит ли дислокаций для быстрой деформации?]
При больших напряжениях подвижны почти все дислокации; плотность в
сильно деформированном металле $\rho\sim10^{11}$~см$^{-2}$. Возьмём
$b=3\cdot10^{-8}$~см и $v\sim10^5$~см/с (меньше скорости звука
$c\sim3\cdot10^5$~см/с):
\[
\dot\varepsilon=\rho bv\approx10^{11}\cdot3\cdot10^{-8}\cdot10^{5}=3\cdot10^{8}~\text{с}^{-1}.
\]
Это означает удвоение длины образца за $\sim3\cdot10^{-9}$~с, тогда как
обычные скорости деформации редко превышают $10^{-4}$--$10^{2}$~с$^{-1}$.
Вывод: дислокаций с запасом хватает для самых быстрых наблюдаемых
деформаций.
\end{primer}

\begin{itogi}
\begin{itemize}
\item Скольжение~— консервативное движение в плоскости, содержащей $\vb b$;
не требует диффузии и идёт при любой $T$. Сила на единицу длины $F=\tau b$.
\item Решётка сопротивляется скольжению: напряжение Пайерлса
$\tau_\text{П}\approx(2G/K)\,e^{-2\pi w/b}$ мало для широких дислокаций,
плотноупакованных плоскостей и коротких $\vb b$. Барьер обходится
зарождением и расхождением двойных перегибов.
\item Переползание~— движение краевой дислокации поперёк плоскости
скольжения за счёт притока/оттока вакансий (через пороги); диффузионное,
термически активируемое. Винтовая дислокация не переползает.
\item Винтовая дислокация может сменить плоскость поперечным скольжением.
\item Формула Орована: $\varepsilon=\rho bl$, $\dot\varepsilon=\rho bv$.
\end{itemize}
\end{itogi}

\begin{proverka}
\item Почему скольжение возможно при любой температуре, а переползание~—
нет?
\item Почему скольжение идёт по плотноупакованным плоскостям?
\item Почему винтовая дислокация не может переползать?
\item В каком направлении движется винтовая дислокация относительно
вектора Бюргерса и в каком направлении смещаются атомы её ядра?
\item Плотность подвижных дислокаций $10^{8}$~см$^{-2}$, $b=2{,}5\cdot10^{-8}$~см.
С какой средней скоростью должны двигаться дислокации, чтобы скорость
деформации была $10^{-3}$~с$^{-1}$?
\end{proverka}
